\SourceLeaf{280}
\begin{list}{}{\leftmargin=2em\rightmargin=0pt\itemindent=0pt\itemsep=0pt\parsep=0pt}
\item first, the triangles are not too far from the standard parallel, and
\item second, they are small in the sense that their sides are short compared with an entire quadrant of the Earth, as is always the case with triangles that can actually be measured.
\end{list}
\Needspace{8\baselineskip}
\noindent\hspace*{\parindent}The triangle sides mapped onto the sphere fit great-circle arcs even more closely in the conformal representation considered above than in the one proposed in the work just cited. At a distance of \(2\tfrac{1}{2}\) degrees from the standard parallel, that earlier representation [according to Art. 13] gave a linear enlargement of \(\tfrac{1}{530000}\), whereas the new method would give a change of only \(\tfrac{1}{5800000}\).

We may therefore first transfer one side of a triangle properly to the surface of the sphere, then use the angles to calculate the whole system just as if it lay on that surface itself, with the modification just indicated if necessary. We can then find the latitudes and longitudes of all the points and pass from these to the latitudes and longitudes on the ellipsoid by means of the formulas given above, or rather, for the latitudes, by means of an auxiliary table such as the one appended here.

\Article{12}
\indent It remains to develop a way of determining how far a geodesic mapped onto the sphere deviates from the great-circle arc between the same endpoints. This also lets us judge in each case whether the deviation needs to be taken into account. There is more than one way to treat this problem. For the present purpose, where the correction can only ever be very small, the following method seems the most suitable.

Let \(L\) be the geodesic in question on the ellipsoid, considered with no fixed limit to its extent, \(M\) its conformal image on the sphere, \(F\) and \(G\) the endpoints of a particular portion of \(M\), and \(N\) a great circle through those two points. Determine each point on \(N\) by its distance \(x\) from an origin chosen arbitrarily on \(N\), and each point on \(M\) by its perpendicular distance \(y\) from \(N\) and by the coordinate of the foot of that perpendicular,
